\documentclass[11pt, a4paper]{article}

% Gemeinsame Style-Definitionen (TEXINPUTS muss templates/ enthalten — siehe scripts/build_tex.py)
% bfw-style.tex — gemeinsame Style-Definitionen für alle Handouts/Aufgabenhefte
% Voraussetzung: Kompilation mit xelatex (wegen fontspec)

\usepackage[ngerman]{babel}
% Babel-ngerman macht " zu einem aktiven Shorthand (z.B. für "a -> ä). In unseren
% Texten verwenden wir " als normales Zeichen — daher Shorthand komplett ausschalten.
\shorthandoff{"}
\usepackage{geometry}
\geometry{a4paper, top=2.4cm, bottom=2.4cm, left=2.3cm, right=2.3cm,
          headheight=15pt, headsep=0.7cm, footskip=1.1cm}

\usepackage{fontspec}
% Fallback-Kette: Source Sans Pro -> Calibri -> Segoe UI -> Latin Modern Sans
% (Latin Modern ist immer mit MiKTeX vorhanden)
\IfFontExistsTF{Source Sans Pro}{%
  \setmainfont{Source Sans Pro}[Ligatures=TeX]%
}{\IfFontExistsTF{Calibri}{%
  \setmainfont{Calibri}[Ligatures=TeX]%
}{\IfFontExistsTF{Segoe UI}{%
  \setmainfont{Segoe UI}[Ligatures=TeX]%
}{%
  \setmainfont{Latin Modern Sans}[Ligatures=TeX]%
}}}
\IfFontExistsTF{Source Code Pro}{%
  \setmonofont{Source Code Pro}[Scale=0.92]%
}{\IfFontExistsTF{Consolas}{%
  \setmonofont{Consolas}[Scale=0.92]%
}{%
  \setmonofont{Latin Modern Mono}[Scale=0.92]%
}}

% Gesperrte Versalien für Labels/Titelbalken (Kapitälchen-Ersatz, funktioniert
% mit jeder OpenType-Schrift — Latin Modern Sans hat keine echten Kapitälchen)
\newcommand{\bfwlabelfont}{\footnotesize\bfseries\addfontfeatures{LetterSpace=8.0}}

\usepackage{xcolor}
% Farbpalette: hell, freundlich, modern
\definecolor{bfwPrimary}{HTML}{0EA5A4}    % Petrol/Türkis - Primär
\definecolor{bfwPrimaryDark}{HTML}{0F766E} % dunkler Petrol für Headings
\definecolor{bfwAccent}{HTML}{F59E0B}     % Amber - Akzent
\definecolor{bfwSuccess}{HTML}{10B981}    % Grün - Erfolg / Lösung
\definecolor{bfwWarn}{HTML}{F97316}       % Orange - Achtung
\definecolor{bfwInk}{HTML}{0F172A}        % fast Schwarz
\definecolor{bfwInkSoft}{HTML}{475569}    % gedämpftes Grau
\definecolor{bfwBgSoft}{HTML}{F8FAFC}     % off-white
\definecolor{bfwBgTeal}{HTML}{ECFEFF}     % helles Türkis
\definecolor{bfwBgAmber}{HTML}{FEF3C7}    % helles Gelb
\definecolor{bfwBgRose}{HTML}{FFE4E6}     % helles Rosé für Eselsbrücken
\definecolor{bfwTabHead}{HTML}{E4F3F2}    % zarte Kopfzeilen-Hinterlegung für Tabellen

\usepackage[colorlinks=true, linkcolor=bfwPrimaryDark, urlcolor=bfwPrimaryDark]{hyperref}
\usepackage{lastpage}
\usepackage{enumitem}
\setlist[itemize]{leftmargin=*, itemsep=2pt, topsep=2pt}
\setlist[enumerate]{leftmargin=*, itemsep=2pt, topsep=2pt}

\usepackage[most]{tcolorbox}
\usepackage{booktabs}
\usepackage{colortbl}
\usepackage{tikz}
\usetikzlibrary{decorations.pathmorphing, positioning, shapes.geometric, arrows.meta}

% ---------------------------------------------------------------------------
% Einheitliches Tabellen-Design
%   - etwas mehr Zeilenluft, dezente graue Linien (wirkt auch mit \hline ruhiger)
%   - Kopfzeile einer Tabelle bei Bedarf zart hinterlegen: \rowcolor{bfwTabHead}
% ---------------------------------------------------------------------------
\renewcommand{\arraystretch}{1.25}
\arrayrulecolor{bfwInkSoft!50}
\setlength{\heavyrulewidth}{1pt}
\setlength{\lightrulewidth}{0.5pt}

% ---------------------------------------------------------------------------
% Boxen — klare farbige Titelbalken mit gesperrten Versal-Labels
% (Titel bleibt per Option überschreibbar: \begin{merksatz}[title={...}])
% ---------------------------------------------------------------------------
\tcbset{bfwbox/.style={%
  enhanced, breakable,
  boxrule=0.5pt, arc=2.5pt,
  left=10pt, right=10pt, top=8pt, bottom=8pt,
  toptitle=4pt, bottomtitle=4pt,
  titlerule=0pt,
  fonttitle=\bfwlabelfont,
  coltitle=white
}}

% Merkkästchen — wichtige Definition / Kernpunkt
\newtcolorbox{merksatz}[1][]{%
  bfwbox,
  colback=bfwBgTeal,
  colframe=bfwPrimaryDark,
  colbacktitle=bfwPrimaryDark,
  title={MERKE},
  #1
}

% Eselsbrücke — Mnemoniks
\newtcolorbox{eselsbruecke}[1][]{%
  bfwbox,
  colback=bfwBgRose,
  colframe=bfwAccent!85!black,
  colbacktitle=bfwAccent!85!black,
  title={ESELSBRÜCKE},
  #1
}

% Beispiel-Box
\newtcolorbox{beispiel}[1][]{%
  bfwbox,
  colback=bfwBgSoft,
  colframe=bfwInkSoft,
  colbacktitle=bfwInkSoft,
  title={BEISPIEL},
  #1
}

% Lösungs-Box (für Aufgabenhefte)
\newtcolorbox{loesung}[1][]{%
  bfwbox,
  colback=bfwSuccess!7,
  colframe=bfwSuccess!65!black,
  colbacktitle=bfwSuccess!65!black,
  title={LÖSUNG},
  #1
}

% Achtung-Box — typische Fehler / Stolperfallen
\newtcolorbox{achtung}[1][]{%
  bfwbox,
  colback=bfwWarn!6,
  colframe=bfwWarn!85!black,
  colbacktitle=bfwWarn!85!black,
  title={ACHTUNG},
  #1
}

% Formel-Box — für zentrale Formeln (groß, ruhig, ohne Titelbalken)
\newtcolorbox{formelbox}[1][]{%
  enhanced,
  colback=bfwBgTeal,
  colframe=bfwPrimaryDark,
  boxrule=1pt, arc=3pt,
  left=10pt, right=10pt, top=6pt, bottom=6pt,
  #1
}

% Glossar-Eintrag
\newcommand{\glossareintrag}[2]{%
  \par\noindent\textbf{\color{bfwPrimaryDark}#1} \enspace #2\par\smallskip
}

% ---------------------------------------------------------------------------
% Abschnitts-Design: farbiges Nummern-Quadrat + feine Linie unter \section,
% dezent farbige \subsection/\subsubsection
% ---------------------------------------------------------------------------
\usepackage{titlesec}
\newcommand{\bfwSecBadge}[1]{%
  \begin{tikzpicture}[baseline={([yshift=-0.62em]n.center)}]
    \node[fill=bfwPrimaryDark, text=white, font=\normalsize\bfseries,
          minimum width=1.55em, minimum height=1.55em, inner sep=1pt,
          rounded corners=1.5pt] (n) {#1};
  \end{tikzpicture}}
\titleformat{\section}
  {\Large\bfseries\color{bfwPrimaryDark}}
  {\bfwSecBadge{\thesection}}{0.6em}{}
  [\vspace{-0.2em}{\color{bfwPrimary!60}\titlerule[0.8pt]}]
\titleformat{\subsection}
  {\large\bfseries\color{bfwPrimaryDark}}
  {\textcolor{bfwPrimary}{\thesubsection}}{0.5em}{}
\titleformat{\subsubsection}
  {\normalsize\bfseries\color{bfwInk}}
  {\textcolor{bfwPrimary}{\thesubsubsection}}{0.4em}{}
\titlespacing*{\section}{0pt}{1.6em}{1em}
\titlespacing*{\subsection}{0pt}{1.2em}{0.5em}

% TikZ-Stil "skizzenhaft" — etwas weicher als Rough.js, aber stabil
\tikzset{
  roughbox/.style = {
    draw=bfwPrimaryDark, thick, fill=bfwBgTeal,
    rounded corners=4pt, inner sep=6pt
  },
  roughaccent/.style = {
    draw=bfwAccent, thick, fill=bfwBgAmber,
    rounded corners=4pt, inner sep=6pt
  },
  roughline/.style = {
    draw=bfwInkSoft, thick, -{Stealth[length=2.5mm]}
  }
}

% Bullet-Symbole (bewusst kein FontAwesome — vermeidet Font-Installations-Probleme)
\newcommand{\faLightbulb}{$\bullet$}
\newcommand{\faBookmark}{$\bullet$}

% ---------------------------------------------------------------------------
% Kopf-/Fußzeile
%   Kopf links:  Wortlogo „Wirtschaftsfachwirt"
%   Kopf rechts: Dokument-Kurztext, gesetzt via \kopfzeile{...}
%   Fuß links:   © Can Siebert · 2026 — Fuß rechts: Seite X von Y
%   Titelseite (Seite 1) bleibt ohne Kopf-/Fußzeile (\handouttitel setzt
%   \thispagestyle{empty}).
% ---------------------------------------------------------------------------
\usepackage{fancyhdr}
\newcommand{\bfwKopfzeileText}{}
\newcommand{\kopfzeile}[1]{\renewcommand{\bfwKopfzeileText}{#1}}
\pagestyle{fancy}
\fancyhf{}
\fancyhead[L]{\small\bfseries\color{bfwPrimaryDark}Wirtschaftsfachwirt}
\fancyhead[R]{\small\color{bfwInkSoft}\bfwKopfzeileText}
\renewcommand{\headrulewidth}{0.6pt}
\renewcommand{\headrule}{{\color{bfwPrimary}\hrule width\headwidth height 0.6pt}}
\fancyfoot[L]{\small\color{bfwInkSoft}\copyright\ Can Siebert\,\textperiodcentered\,2026}
\fancyfoot[R]{\small\color{bfwInkSoft}Seite \thepage\ von \pageref*{LastPage}}
% Auch \thispagestyle{plain} (z.B. durch \maketitle) soll leer bleiben:
\fancypagestyle{plain}{\fancyhf{}\renewcommand{\headrulewidth}{0pt}\renewcommand{\headrule}{}}

% ---------------------------------------------------------------------------
% \handouttitel{Obertitel}{Untertitel}{Kurs-Label}{Termin-Zeile}
%   Einheitlicher professioneller Titelblock: Dachzeile, farbige Headline,
%   Untertitel, farbige Trennlinie und dezente Meta-Zeile (Kurs/Termin/Dozent).
%   Ersetzt \title/\maketitle. Direkt nach \begin{document} aufrufen.
% ---------------------------------------------------------------------------
\newcommand{\handouttitel}[4]{%
  \thispagestyle{empty}%
  \begingroup
  \setlength{\parindent}{0pt}%
  {\bfwlabelfont\color{bfwPrimary}WIRTSCHAFTSFACHWIRT (IHK)\par}
  \vspace{0.55em}
  {\huge\bfseries\color{bfwPrimaryDark}#1\par}
  \vspace{0.5em}
  {\large\color{bfwInkSoft}#2\par}
  \vspace{0.9em}
  {\color{bfwPrimary}\rule{\linewidth}{2pt}}\par
  \vspace{0.55em}
  {\small\color{bfwInkSoft}#3\ \,\textperiodcentered\,\ #4\ \,\textperiodcentered\,\ Dozent: Can Siebert\par}
  \endgroup
  \vspace{1.4em}%
}


\title{\vspace*{-1.5cm}\Huge\bfseries\color{bfwPrimaryDark}Lektion~4: Unternehmenszusammenschlüsse\\[0.4em]
  \large\color{bfwInkSoft}\textnormal{Von der Kooperation zur Konzentration --- und wo das Kartellrecht Grenzen zieht}}
\author{\textsf{Wirtschaftsfachwirt --- Betriebswirtschaftliche Grundlagen \textperiodcentered{} \small\copyright\ Can Siebert\,\textperiodcentered\,2027}}
\date{}

\begin{document}
\maketitle
\thispagestyle{empty}

\vspace{1em}
\begin{merksatz}[title={\bfseries Worum geht es in dieser Lektion?}]
Unternehmen wachsen nicht nur aus eigener Kraft (internes Wachstum), sondern auch durch
Verbindungen mit anderen Unternehmen (externes Wachstum). In dieser Lektion lernen Sie,
warum sich Unternehmen zusammenschließen, welche Formen es gibt und woran man sie
unterscheidet: Bei der \textbf{Kooperation} bleiben die Partner selbstständig, bei der
\textbf{Konzentration} geht die wirtschaftliche --- teils auch die rechtliche ---
Selbstständigkeit verloren. Außerdem ordnen Sie Zusammenschlüsse nach ihrer Richtung
(horizontal, vertikal, diagonal) ein und lernen die Grenzen kennen, die Kartellverbot und
Fusionskontrolle setzen. Die Lektion baut unmittelbar auf den Rechtsformen (Lektion~3)
auf und schließt den Rahmenplan-Abschnitt 1.4 ab.
\end{merksatz}

\begin{beispiel}[title={Hinweis zur Literatur}]
Der Textband-Auszug zu diesem Thema wird nachgereicht. Bis dahin dient dieses Handout als
Skriptersatz; ergänzend: Übungsband Betriebswirtschaft, Übungen 54--58 (mit Lösungen).
\end{beispiel}

\tableofcontents
\vspace{1em}
\hrule
\vspace{1em}

\section{Anknüpfung: Rechtsformen als Bausteine (Wiederholung Lektion 3)}

Zusammenschlüsse verbinden \emph{bestehende} Unternehmen --- also die Rechtsformen aus
Lektion~3. Zur Erinnerung:

\begin{itemize}
  \item \textbf{Personengesellschaften} (GbR, OHG, KG): keine juristischen Personen,
    persönliche Haftung mindestens eines Gesellschafters. Die \emph{GbR} ist die typische
    Rechtsform für Arbeitsgemeinschaften und Konsortien.
  \item \textbf{Kapitalgesellschaften} (GmbH, AG): juristische Personen, Haftung auf das
    Gesellschaftsvermögen beschränkt. Sie sind die typischen Bausteine von
    \emph{Konzernen} (Mutter- und Tochtergesellschaften) und von \emph{Joint Ventures}
    (Gemeinschafts-GmbH).
  \item Die \textbf{AG} kann Anteile anderer Gesellschaften erwerben ---
    Mehrheitsbeteiligungen sind der häufigste Weg in den Konzern.
\end{itemize}

\begin{merksatz}
Ein Zusammenschluss ändert nicht automatisch die Rechtsform: Im \emph{Konzern} bleiben
alle Gesellschaften rechtlich bestehen; erst bei der \emph{Fusion (Verschmelzung)}
erlischt mindestens ein Rechtsträger.
\end{merksatz}

\section{Gründe und Ziele von Unternehmenszusammenschlüssen}

Unternehmen schließen sich zusammen, um gemeinsam Ziele zu erreichen, die allein nicht
oder nur schwer erreichbar wären:

\begin{itemize}
  \item \textbf{Rationalisierung und Kostensenkung:} Größenvorteile (Kostendegression),
    gemeinsamer Einkauf (Mengenrabatte), gemeinsame Verwaltung, Vertrieb, Logistik;
  \item \textbf{Synergieeffekte:} Zusammenwirken schafft Mehrwert („2 + 2 = 5``) ---
    z.\,B. Teilung der Forschungs- und Entwicklungskosten, Wissens- und
    Technologietransfer;
  \item \textbf{Marktmacht und Wettbewerbsposition:} größere Marktanteile, Ausschalten
    von Konkurrenz, stärkere Position gegenüber Lieferanten und Abnehmern;
  \item \textbf{Sicherung von Beschaffung und Absatz:} Zugang zu Rohstoffen,
    Vorprodukten und Vertriebswegen (vertikale Motive);
  \item \textbf{Risikostreuung:} Verteilung des unternehmerischen Risikos auf mehrere
    Geschäftsfelder oder Projekte (Diversifikation);
  \item \textbf{Kapitalkraft und Finanzierung:} bessere Kreditkonditionen, gemeinsame
    Finanzierung von Großprojekten;
  \item \textbf{Schutz vor Übernahmen} und Erschließung neuer (Auslands-)Märkte.
\end{itemize}

\begin{eselsbruecke}
\textbf{SMaRRt} zusammenschließen: \textbf{S}ynergien, \textbf{Ma}rktmacht,
\textbf{R}ationalisierung, \textbf{R}isikostreuung --- die vier Kernziele, die die IHK
am liebsten abfragt.
\end{eselsbruecke}

\section{Das Schlüsselkriterium: Selbstständigkeit}

Alle Zusammenschlussformen werden danach eingeteilt, was mit der
\textbf{rechtlichen} und der \textbf{wirtschaftlichen} Selbstständigkeit der beteiligten
Unternehmen geschieht:

\begin{itemize}
  \item \textbf{Rechtliche Selbstständigkeit:} Das Unternehmen bleibt eigener
    Rechtsträger (eigene Firma, eigene Organe, eigene Bilanz).
  \item \textbf{Wirtschaftliche Selbstständigkeit:} Das Unternehmen trifft seine
    unternehmerischen Entscheidungen (Preise, Programme, Investitionen) selbst.
\end{itemize}

\subsection{Kooperation vs. Konzentration im Überblick}

\begin{center}
\begin{tabular}{p{4.2cm} p{4.9cm} p{4.9cm}}
\toprule
\textbf{Kriterium} & \textbf{Kooperation} & \textbf{Konzentration}\\
\midrule
Rechtliche Selbstständigkeit & bleibt erhalten & Konzern: bleibt erhalten;\newline Fusion/Trust: geht (bei mind. einem Partner) verloren\\
\addlinespace[3pt]
Wirtschaftliche Selbstständigkeit & bleibt grundsätzlich erhalten (nur im Kooperationsbereich freiwillig eingeschränkt) & geht verloren (einheitliche Leitung bzw. ein einziges Unternehmen)\\
\addlinespace[3pt]
Bindungsintensität & locker bis vertraglich, meist befristet/kündbar & dauerhaft, kapital- oder vertragsbasiert\\
\addlinespace[3pt]
Typische Formen & Arbeitsgemeinschaft/Konsortium, Interessengemeinschaft, Joint Venture, strategische Allianz, Verband, Kartell & Konzern, Fusion (Verschmelzung), Trust\\
\addlinespace[3pt]
Typische Ziele & Synergien nutzen, Kosten teilen, Projekte stemmen --- bei voller Eigenständigkeit & Marktmacht, einheitliche Leitung, volle Kontrolle, Rationalisierung\\
\addlinespace[3pt]
Kartellrechtliche Sicht & grundsätzlich erlaubt --- \emph{außer} wettbewerbsbeschränkende Absprachen (Kartellverbot) & ab bestimmten Größen anmeldepflichtig: Fusionskontrolle\\
\bottomrule
\end{tabular}
\end{center}

\begin{merksatz}[title={\bfseries Merksatz: Die Kernabgrenzung}]
\textbf{Kooperation} = freiwillige Zusammen\emph{arbeit} rechtlich und wirtschaftlich
selbstständiger Unternehmen. \textbf{Konzentration} = Zusammen\emph{fassung} von
Unternehmen unter Verlust der wirtschaftlichen Selbstständigkeit --- beim Konzern bleibt
die rechtliche Selbstständigkeit bestehen, bei Fusion und Trust geht auch sie verloren.
\end{merksatz}

\section{Formen der Kooperation}

\subsection{Arbeitsgemeinschaft (ARGE) und Konsortium}

Rechtlich und wirtschaftlich selbstständige Unternehmen schließen sich \emph{befristet}
zusammen, um \textbf{einen einzelnen Auftrag oder ein Projekt} gemeinsam durchzuführen.
Rechtsform ist regelmäßig die \textbf{GbR}; nach Projektende wird der Zusammenschluss
aufgelöst.

\begin{itemize}
  \item \textbf{Arbeitsgemeinschaft:} typisch in der \emph{Bauwirtschaft} --- mehrere
    Bauunternehmen errichten gemeinsam einen Tunnel oder eine Brücke, weil Kapazität,
    Know-how und Risiko für ein einzelnes Unternehmen zu groß wären.
  \item \textbf{Konsortium:} derselbe Grundgedanke im \emph{Bank- und Finanzwesen} ---
    z.\,B. ein Bankenkonsortium zur gemeinsamen Emission einer Anleihe oder zur
    Finanzierung eines Großkredits.
\end{itemize}

\subsection{Interessengemeinschaft (IG)}

Vertragliche, auf \emph{Dauer} angelegte Zusammenarbeit rechtlich selbstständiger
Unternehmen in \textbf{einzelnen Teilbereichen} --- z.\,B. gemeinsame Forschung und
Entwicklung, gemeinsame Patentverwertung oder ein Gemeinschaftslabor. Häufig wird ein
Gewinn- oder Kostenpool für den gemeinsamen Bereich vereinbart. Die Unternehmen bleiben
im Übrigen unabhängig und treten am Markt weiter getrennt auf.

\subsection{Joint Venture (Gemeinschaftsunternehmen)}

Zwei (oder mehr) Partnerunternehmen gründen ein \textbf{rechtlich selbstständiges
Gemeinschaftsunternehmen}, an dem sie sich mit Kapital beteiligen und dessen Führung und
Risiko sie gemeinsam tragen (oft 50:50). Typische Motive: Erschließung von
\emph{Auslandsmärkten} (Partner bringt Marktkenntnis und Vertriebsnetz ein), Teilung hoher
Entwicklungskosten, Kombination komplementärer Stärken.

\begin{beispiel}
Ein deutscher Automobilhersteller und ein chinesischer Partner gründen gemeinsam eine
Produktions- und Vertriebsgesellschaft in China: Beide Mütter bleiben selbstständig ---
das Joint Venture selbst ist eine neue, eigene Gesellschaft.
\end{beispiel}

\subsection{Strategische Allianz}

Zusammenarbeit --- häufig zwischen (potenziellen) \emph{Konkurrenten} --- in
\textbf{einzelnen strategischen Geschäftsfeldern}, ohne gemeinsames
Tochterunternehmen: z.\,B. Luftfahrt-Allianzen (gemeinsames Streckennetz, Codesharing,
Vielfliegerprogramme) oder Entwicklungs-Partnerschaften bei Batteriezellen. Ziel ist es,
Stärken zu bündeln und Wettbewerbsvorteile gegenüber Dritten zu erlangen; außerhalb der
Allianz bleiben die Partner Wettbewerber.

\subsection{Verband}

Freiwilliger Zusammenschluss zur \textbf{gemeinsamen Interessenvertretung nach außen} ---
gegenüber Staat, Öffentlichkeit, Tarifpartnern --- sowie zur Beratung und Information der
Mitglieder. Beispiele: Branchenverbände (z.\,B. Automobilindustrie), Arbeitgeberverbände,
Industrie- und Handelskammern als Körperschaften. Verbände koordinieren \emph{nicht} das
Marktverhalten ihrer Mitglieder --- sonst drohte ein Kartellverstoß.

\subsection{Kartell --- Kooperation mit kartellrechtlicher Sprengkraft}

Ein \textbf{Kartell} ist eine \emph{vertragliche Absprache} (oder abgestimmte
Verhaltensweise) rechtlich selbstständiger Unternehmen derselben Wirtschaftsstufe mit dem
Ziel, den \textbf{Wettbewerb zu beschränken} --- z.\,B.:

\begin{itemize}
  \item \textbf{Preiskartell:} einheitliche Preise oder Preisuntergrenzen,
  \item \textbf{Quoten-/Produktionskartell:} Aufteilung von Produktionsmengen,
  \item \textbf{Gebietskartell:} Aufteilung von Absatzgebieten oder Kunden,
  \item \textbf{Submissionsabsprachen:} Absprachen bei Ausschreibungen.
\end{itemize}

\begin{merksatz}[title={\bfseries Kartellrechtliche Einordnung (Verweis: VWL-Lektion 2)}]
Nach \textbf{\S\,1 GWB} (Gesetz gegen Wettbewerbsbeschränkungen) und
\textbf{Art.\,101 AEUV} sind wettbewerbsbeschränkende Vereinbarungen
\textbf{grundsätzlich verboten} (Verbotsprinzip). Ausnahmen (Freistellung, \S\,2 GWB)
gelten nur, wenn die Vorteile überwiegen --- etwa bei bestimmten
\emph{Mittelstandskooperationen}, die die Wettbewerbsfähigkeit kleiner und mittlerer
Unternehmen verbessern und den Wettbewerb nicht wesentlich beeinträchtigen. Das
\textbf{Bundeskartellamt} verfolgt Verstöße mit empfindlichen Bußgeldern
(Kronzeugenregelung für Aufdecker). Arbeitsgemeinschaften und Konsortien sind dagegen
\emph{grundsätzlich erlaubt}, weil sie den Wettbewerb regelmäßig nicht spürbar
beschränken.
\end{merksatz}

\section{Formen der Konzentration}

\subsection{Konzern}

Ein \textbf{Konzern} ist die Zusammenfassung eines \emph{herrschenden} und eines oder
mehrerer \emph{abhängiger} Unternehmen unter \textbf{einheitlicher Leitung}
(\S\,18 AktG). Die Konzernunternehmen bleiben \textbf{rechtlich selbstständig}, verlieren
aber ihre \textbf{wirtschaftliche Selbstständigkeit}.

\begin{itemize}
  \item \textbf{Muttergesellschaft (Obergesellschaft):} übt die einheitliche Leitung aus
    --- meist über eine \emph{Mehrheitsbeteiligung} an den Töchtern; eine reine
    \emph{Holding} beschränkt sich auf das Halten und Steuern der Beteiligungen.
  \item \textbf{Tochtergesellschaften:} rechtlich selbstständige Gesellschaften
    (typisch GmbH oder AG), die den Weisungen der Mutter folgen.
  \item \textbf{Beherrschungsvertrag} (\S\,291 AktG): Unternehmensvertrag, mit dem sich
    eine Gesellschaft der Leitung eines anderen Unternehmens unterstellt --- die Mutter
    erhält ein direktes \emph{Weisungsrecht} gegenüber dem Vorstand der Tochter; häufig
    kombiniert mit einem Gewinnabführungsvertrag.
  \item \textbf{Unterordnungskonzern} (Regelfall: Mutter--Tochter) vs.
    \textbf{Gleichordnungskonzern} (Unternehmen unter einheitlicher Leitung, ohne dass
    eines vom anderen abhängig ist).
\end{itemize}

\subsection{Fusion (Verschmelzung)}

Die \textbf{Fusion} ist die engste Form des Zusammenschlusses: Mindestens ein beteiligtes
Unternehmen gibt seine \textbf{rechtliche und wirtschaftliche Selbstständigkeit} auf
(geregelt im Umwandlungsgesetz). Zwei Wege:

\begin{itemize}
  \item \textbf{Verschmelzung durch Aufnahme:} Unternehmen B überträgt sein Vermögen auf
    Unternehmen A; B erlischt (A + B $\rightarrow$ A).
  \item \textbf{Verschmelzung durch Neugründung:} A und B übertragen ihr Vermögen auf
    eine neu gegründete Gesellschaft C; A und B erlöschen (A + B $\rightarrow$ C).
\end{itemize}

\subsection{Trust}

Historischer, v.\,a. aus den USA stammender Begriff für die \textbf{vollständige
Verschmelzung} ehemals selbstständiger Unternehmen zu einer wirtschaftlichen und
rechtlichen Einheit mit marktbeherrschender Stellung. Die beteiligten Unternehmen
verlieren rechtliche \emph{und} wirtschaftliche Selbstständigkeit; Ziel ist die
Monopolisierung ganzer Märkte --- daher der Begriff „Antitrust-Recht`` für das
Wettbewerbsrecht.

\begin{eselsbruecke}
Intensitäts-Treppe der Konzentration: \textbf{K}onzern (rechtlich noch selbstständig)
$\rightarrow$ \textbf{F}usion (ein Unternehmen erlischt) $\rightarrow$ \textbf{T}rust
(alles verschmilzt zur beherrschenden Einheit): \textbf{K}ommt \textbf{f}ast
\textbf{t}otal zusammen.
\end{eselsbruecke}

\section{Richtungen des Zusammenschlusses}

Unabhängig von der Form (Kooperation oder Konzentration) unterscheidet man, \emph{welche}
Unternehmen sich verbinden:

\begin{itemize}
  \item \textbf{Horizontal:} Unternehmen der \emph{gleichen} Branche und
    \emph{gleichen} Produktions- oder Handelsstufe --- z.\,B. zwei Brauereien, zwei
    Großküchenhersteller. Ziele: Marktanteil, Größenvorteile, Rationalisierung.
  \item \textbf{Vertikal:} Unternehmen \emph{aufeinanderfolgender} Wirtschaftsstufen ---
    \emph{vorgelagert} (Richtung Rohstoff/Zulieferer, „Rückwärtsintegration``: Brauerei +
    Mälzerei) oder \emph{nachgelagert} (Richtung Absatz/Kunde, „Vorwärtsintegration``:
    Brauerei + Getränkegroßhandel). Ziele: Sicherung von Beschaffung und Absatz.
  \item \textbf{Diagonal} (auch \emph{anorganisch}, \emph{lateral},
    \emph{konglomerat}): \emph{branchenfremde} Unternehmen ohne
    Leistungszusammenhang --- z.\,B. Maschinenbauer + Versicherung. Ziele:
    Risikostreuung (Diversifikation), Kapitalanlage, neue Geschäftsfelder.
\end{itemize}

\begin{beispiel}[title={Beispiel aus dem Übungsband (Übung 56): Großküchenhersteller}]
Ein Großküchenhersteller verbindet sich \emph{vertikal vorgelagert} mit seinem
Spülmaschinen-Zulieferer, \emph{vertikal nachgelagert} mit einem Kantinenbetreiber
(Abnehmer), \emph{horizontal} mit einem anderen Großküchenhersteller und
\emph{diagonal} mit einer Geschäftsbank (branchenfremd).
\end{beispiel}

\section{Ziele im Vergleich und Grenzen: die Fusionskontrolle}

\subsection{Ziele der Kooperation vs. Ziele der Konzentration}

\begin{center}
\begin{tabular}{p{6.8cm} p{6.8cm}}
\toprule
\textbf{Kooperation --- „gemeinsam stärker, aber frei``} & \textbf{Konzentration --- „Einheit und Kontrolle``}\\
\midrule
Synergien in Teilbereichen nutzen (F\&E, Einkauf, Vertrieb) & umfassende Rationalisierung unter einheitlicher Leitung\\
Großprojekte gemeinsam stemmen (ARGE, Konsortium) & Marktmacht durch Marktanteilsaddition\\
Risiko einzelner Projekte teilen & Risikostreuung durch Diversifikation (konglomerat)\\
Märkte gemeinsam erschließen (Joint Venture, Allianz) & Sicherung von Beschaffung und Absatz durch Integration\\
Selbstständigkeit und Flexibilität bleiben erhalten & volle Kontrolle, aber Verlust der Eigenständigkeit\\
\bottomrule
\end{tabular}
\end{center}

\subsection{Fusionskontrolle durch Bundeskartellamt und EU-Kommission}

Konzentration kann den Wettbewerb dauerhaft ausschalten --- deshalb unterliegen größere
Zusammenschlüsse der \textbf{präventiven Fusionskontrolle}:

\begin{itemize}
  \item \textbf{Deutschland} (\S\S\,35\,ff. GWB): Zusammenschlüsse müssen beim
    \textbf{Bundeskartellamt} (Bonn) \emph{vor} dem Vollzug \emph{angemeldet} werden,
    wenn die beteiligten Unternehmen bestimmte Umsatzschwellen überschreiten (weltweit
    zusammen mehr als 500~Mio.~€; Inlandsumsätze der Beteiligten über weiteren
    Schwellenwerten).
  \item \textbf{Untersagung:} Der Zusammenschluss wird verboten, wenn er wirksamen
    Wettbewerb \emph{erheblich behindern} würde --- insbesondere wenn eine
    \textbf{marktbeherrschende Stellung} entsteht oder verstärkt wird (Vermutung ab
    40\,\% Marktanteil eines Unternehmens). Möglich sind auch Freigaben unter
    \emph{Auflagen} (z.\,B. Verkauf von Unternehmensteilen).
  \item \textbf{Ministererlaubnis} (\S\,42 GWB): Ausnahmsweise kann das
    Bundeswirtschaftsministerium einen untersagten Zusammenschluss aus
    Gemeinwohlgründen dennoch erlauben.
  \item \textbf{EU-Ebene:} Bei \emph{gemeinschaftsweiter Bedeutung} (sehr hohe,
    europaweite Umsätze) prüft die \textbf{EU-Kommission} nach der
    EU-Fusionskontrollverordnung (FKVO) --- dann ist sie anstelle des Bundeskartellamts
    zuständig („One-Stop-Shop``).
\end{itemize}

\begin{merksatz}[title={\bfseries Wichtig für die Prüfungsvorbereitung}]
Sie sollten sicher können: die Abgrenzung \emph{Kooperation vs. Konzentration} über die
rechtliche und wirtschaftliche Selbstständigkeit, je drei \emph{Formen mit Beispiel},
die Definition des \emph{Konzerns} (rechtlich selbstständig, einheitliche Leitung,
Beherrschungsvertrag), die \emph{Fusionswege} (Aufnahme/Neugründung), die drei
\emph{Richtungen} mit Praxisbeispielen sowie \emph{Kartellverbot} (\S\,1 GWB) und
\emph{Fusionskontrolle} (Bundeskartellamt/EU-Kommission). Üben Sie mit dem Aufgabenheft,
dem Quiz und den Übungen 54--58 des Übungsbands!
\end{merksatz}

\vspace{1em}
\hrule
\vspace{0.8em}
\noindent\textsf{\small Quellen: Rahmenplan Wirtschaftsbezogene Qualifikationen (DIHK), Abschnitt 1.4;
Übungsband Betriebswirtschaft (DIHK-Bildungs-GmbH), Übungen 54--58; GWB, AktG, UmwG (Grundzüge).
Textband-Auszug wird nachgereicht. Dies ist der letzte Kurstermin --- viel Erfolg bei der IHK-Prüfung!}

\end{document}
